\documentclass[11pt,a4paper]{article}

% ---------------------------------------------------------
% PREAMBLE (stable, warning-free, Overleaf-safe)
% ---------------------------------------------------------
\usepackage[utf8]{inputenc}
\usepackage[T1]{fontenc}
\usepackage{lmodern}
\usepackage{amsmath, amssymb, amsfonts}
\usepackage{bm}
\usepackage{geometry}
\usepackage{graphicx}
\usepackage{hyperref}
\usepackage{cite}
\usepackage{authblk}
\usepackage{physics}
\usepackage{mathtools}

\geometry{margin=2.4cm}

\title{\textbf{Companion LXIV: Wasserstein Distances of Persistence Diagrams in the Relaxation-Driven Cyclic Cosmology}}
\author{Michael Lehmann \\ \texttt{mi.lehmann@gmx.de}}
\date{April 2026}

\begin{document}
\maketitle

\begin{abstract}
Wasserstein distances provide a metric-space comparison of persistence diagrams, 
quantifying the difference between topological structures across filtrations. In 
weak-lensing analyses, these distances measure how the topology of convergence 
fields changes with redshift, smoothing scale, or cosmological model. In the 
standard cosmological model, Wasserstein distances reflect the nonlinear matter 
distribution and the projection of large-scale structure. In the Relaxation-Driven 
Cyclic Cosmology (RDCC), the scale-dependent evolution of the gravitational 
potential modifies the metric structure of persistence diagrams in a 
characteristic way. This Companion develops a continuous description of 
Wasserstein distances in RDCC, deriving how the relaxation scale influences 
diagram geometry, topological stability, and the observational signatures in 
weak-lensing surveys. The resulting framework provides a distinctive metric-space 
probe of the RDCC gravitational field.
\end{abstract}
% ---------------------------------------------------------
\section{Introduction}

Persistent homology provides a topological description of cosmological fields by 
tracking the birth and death of connected components and loops across a 
filtration. Wasserstein distances extend this by defining a metric on the space of 
persistence diagrams, allowing a quantitative comparison of topological structure.

In the standard cosmological model, Wasserstein distances reflect the nonlinear 
matter distribution and the projection of large-scale structure. In RDCC, the 
gravitational potential evolves differently. The relaxation mechanism suppresses 
the decay of small-scale modes and enhances the coherence of large-scale modes. 
This modifies the metric structure of persistence diagrams in a scale-dependent 
manner, producing a distinctive signature in Wasserstein-distance analyses.
% ---------------------------------------------------------
\section{Persistence Diagrams and Wasserstein Distances}

A persistence diagram $D$ consists of points $(b,d)$ representing the birth and 
death of topological features. The $p$-Wasserstein distance between two diagrams 
$D_{1}$ and $D_{2}$ is

\begin{equation}
    W_{p}(D_{1},D_{2})
    = \left[
        \inf_{\gamma}
        \sum_{x \in D_{1}}
        \| x - \gamma(x) \|_{\infty}^{p}
      \right]^{1/p},
\end{equation}

where $\gamma$ ranges over all bijections between the diagrams (including the 
diagonal).
% ---------------------------------------------------------
\section{RDCC Modification of Persistence Diagrams}

In RDCC, the convergence evolves as
\begin{equation}
    \kappa_{\mathrm{RDCC}}(\ell)
    = \kappa(\ell)
      \left[
        1 - \chi_{\kappa}
        \left(\frac{\ell}{\ell_{\mathrm{rel}}}\right)^{\alpha}
      \right].
\end{equation}

This modifies the birth and death thresholds of topological features:

\begin{align}
    b_{\mathrm{RDCC}} &= b \left[ 1 - \chi_{b} (\ell/\ell_{\mathrm{rel}})^{\alpha} \right], \\
    d_{\mathrm{RDCC}} &= d \left[ 1 - \chi_{d} (\ell/\ell_{\mathrm{rel}})^{\alpha} \right].
\end{align}

The RDCC persistence diagram becomes
\begin{equation}
    D_{\mathrm{RDCC}}
    = \left\{
        (b_{\mathrm{RDCC}}, d_{\mathrm{RDCC}})
      \right\}.
\end{equation}
% ---------------------------------------------------------
\section{Wasserstein Distances in RDCC}

The Wasserstein distance between a standard and an RDCC diagram is

\begin{equation}
    W_{p}(D, D_{\mathrm{RDCC}})
    = \left[
        \sum_{i}
        \left\|
          (b_{i}, d_{i})
          -
          (b_{i,\mathrm{RDCC}}, d_{i,\mathrm{RDCC}})
        \right\|_{\infty}^{p}
      \right]^{1/p}.
\end{equation}

This becomes
\begin{equation}
    W_{p}(D, D_{\mathrm{RDCC}})
    = W_{p}(D,D)
      \left[
        1 - \chi_{W}
        \left(\frac{\ell}{\ell_{\mathrm{rel}}}\right)^{\alpha}
      \right],
\end{equation}

reflecting the relaxation-driven modification of topological distances.
% ---------------------------------------------------------
\section{Metric-Space Signatures of RDCC}

The relaxation mechanism enhances the coherence of large-scale modes. Wasserstein 
distances therefore exhibit:

\begin{itemize}
    \item reduced distances between diagrams at large scales,
    \item enhanced distances at small scales,
    \item a characteristic turnover near $\ell_{\mathrm{rel}}$,
    \item modified scaling with smoothing scale and redshift,
    \item altered clustering of diagrams in metric space.
\end{itemize}

These signatures provide a direct observational probe of the RDCC gravitational 
field in metric space.
% ---------------------------------------------------------
\section{Conclusion}

Wasserstein distances of persistence diagrams in the Relaxation-Driven Cyclic 
Cosmology reflect the scale-dependent evolution of the gravitational potential and 
the nonlinear matter distribution. The relaxation mechanism modifies the metric 
structure of persistence diagrams in a scale-dependent manner, producing 
distinctive signatures in weak-lensing surveys. These effects provide a direct 
observational window into the relaxation scale and the temporal structure of the 
RDCC gravitational field. The present Companion establishes the theoretical 
foundation for future studies of topological metric spaces, nonlinear structure, 
and the interplay between light and gravity in RDCC.

\bibliographystyle{apsrev4-2}
\nocite{*}
\bibliography{references}


\end{document}
